\documentclass[12pt, a4paper]{article}

% Paquetes esenciales
\usepackage[utf8]{inputenc}
\usepackage[spanish]{babel}
\usepackage{geometry}
\geometry{top=2.5cm, bottom=2.5cm, left=3cm, right=3cm}
\usepackage{graphicx}
\usepackage{hyperref}
\usepackage{titlesec}
\usepackage{xcolor}
\usepackage{enumitem}
\usepackage{parskip}       % Elimina la sangría inicial y añade espacio entre párrafos

% Paquetes agregados para este informe
\usepackage{booktabs}      % tablas con lineas limpias
\usepackage{longtable}     % tablas que pueden partir de pagina
\usepackage{array}         % columnas mejor definidas
\usepackage{listings}      % bloques de codigo
\usepackage{float}         % figuras fijas
\usepackage{caption}       % pie de figura
\usepackage[most]{tcolorbox} % recuadros de aviso
\usepackage{tikz}          % diagramas
\usetikzlibrary{arrows.meta} % puntas de flecha

% Configuración de enlaces
\hypersetup{
    colorlinks=true,
    linkcolor=blue,
    filecolor=magenta,
    urlcolor=blue,
}

% ---- Configuracion de tablas ----
\renewcommand{\arraystretch}{1.25}

% ---- Configuracion de bloques de codigo ----
\definecolor{codebg}{RGB}{246,247,249}
\definecolor{codeborder}{RGB}{210,214,220}
\lstset{
    backgroundcolor=\color{codebg},
    frame=single,
    rulecolor=\color{codeborder},
    basicstyle=\ttfamily\small,
    keywordstyle=\color{blue!70!black}\bfseries,
    commentstyle=\color{green!40!black},
    stringstyle=\color{red!50!black},
    showstringspaces=false,
    breaklines=true,
    columns=fullflexible,
    framexleftmargin=6pt,
    framexrightmargin=6pt,
    xleftmargin=4pt,
    aboveskip=10pt,
    belowskip=10pt,
}

% ---- Recuadro de aviso ----
\newtcolorbox{aviso}{
    colback=yellow!8,
    colframe=orange!60!black,
    fonttitle=\bfseries,
    title={Espacio reservado},
    breakable,
}

% ==========================================
% 1. PORTADA
% ==========================================

\begin{document}

\titlespacing*{\section}{0pt}{2ex plus 0.5ex minus 0.5ex}{4ex plus 0.5ex}
\titlespacing*{\subsection}{0pt}{2ex plus 0.5ex minus 0.5ex}{2ex plus 0.5ex}
\titlespacing*{\subsubsection}{0pt}{2ex plus 0.5ex minus 0.5ex}{0.5ex plus 0.5ex}

\begin{titlepage}

    \begin{flushleft}
        \vspace*{-1.5em}
        % Si el escudo no esta en la carpeta, la portada compila igual.
        \IfFileExists{escudo.png}{%
            \includegraphics[scale=0.2]{escudo.png}%
        }{%
            \vspace{0.6cm}%
        }
    \end{flushleft}

    \vspace*{4cm}
    \centering
    {\Huge \textbf{Integración Sistemas}} \\[5ex]
    {\large \textit{Encargo 1}} \\[0.1cm]
    {\large \textit{Veterinaria (Forma B)}} \\[10cm]
    \large
    Sebastián Ignacio Vega Varela\\
    Cristóbal Alonso González Cifuentes\\
    Braian Alejandro Urra Bastías\\
    Matías Felipe Jenner Valdebenito Valenzuela
    \vspace{1cm}

    \textbf{Código:}\\
    \url{https://github.com/Sebavegs2004/Entrega1-IntegracionSistemas}

\end{titlepage}

\tableofcontents

\newpage

% ==========================================
% 2. INTRODUCCION
% ==========================================

\section{Introducción}

\subsection{Contexto}

La Cadena de Clínicas Veterinarias VidaAnimal maneja hoy dos sistemas
que se crearon por separado y que nunca se comunicaron entre sí. El
primero es \textbf{Agenda}, que administra los veterinarios y lleva la
cuenta de los cupos libres de cada bloque horario. El segundo es
\textbf{Reservas}, que registra a los dueños de mascotas y las horas que
reservan.

Como los dos sistemas nunca hablaron, el proceso de tomar una hora era
manual: alguien llamaba por teléfono a la clínica para preguntar si había
disponibilidad. Cuando la recepción anotaba mal, o cuando dos personas
preguntaban al mismo tiempo, se producían \textbf{sobrecupos}, es decir,
más reservas que cupos disponibles. El problema no era de código sino de
integración: dos bases de datos que guardan información relacionada y
que nadie mantiene sincronizadas.

\subsection{Qué se construyó}

El equipo construyó una capa de integración entre los dos sistemas,
compuesta por tres piezas:

\begin{itemize}[leftmargin=1.4cm]
    \item Una \textbf{API REST pública} del servicio Reservas, que es la
    puerta de entrada para el personal de la clínica y para la interfaz
    web. Está versionada bajo \texttt{/v1} y documentada en OpenAPI.
    \item Un \textbf{servicio interno en gRPC} que corresponde a Agenda, que
    sigue siendo la fuente de verdad sobre los cupos. No se publica hacia
    el exterior: solo se usa dentro de la red de Docker.
    \item La \textbf{integración entre ambos}, que se produce cada vez que
    un cliente pide reservar una hora: primero se consulta el cupo a Agenda
    y solo después se guarda la reserva.
\end{itemize}

Además se agregó una caché con Redis, idempotencia en el registro de
reservas y enlaces de acción (HATEOAS) en las respuestas.

\subsection{Cómo está organizado este documento}

El documento sigue el orden en que se resolvió el problema. La sección
\ref{sec:analisis} describe qué se pedía y con qué restricciones. La
sección \ref{sec:diseno} explica cómo quedó la solución, con el
diagrama de arquitectura y la justificación de la frontera entre los
servicios. La sección \ref{sec:adr} presenta las cuatro decisiones de
diseño como \textit{Architecture Decision Records} (ADR). La sección
\ref{sec:experimento} corresponde a la experimentación. La sección
\ref{sec:contexto} trata el factor de contexto considerado. Las secciones
finales reúnen las instrucciones de ejecución y la reflexión del equipo.

% ==========================================
% 3. ANALISIS DEL PROBLEMA
% ==========================================

\newpage
\section{Análisis del problema}\label{sec:analisis}

\subsection{Situación inicial}

Antes de este trabajo, los dos sistemas se comportaban así:

\begin{itemize}[leftmargin=1.4cm]
    \item \textbf{Agenda} era el único que sabía cuántos cupos quedaban.
    Nadie más lo consultaba de forma automática.
    \item \textbf{Reservas} guardaba las horas acordadas, pero no tenía
    forma de confirmar que quedara cupo antes de anotarlas.
    \item La \textbf{consulta de disponibilidad} era manual: la recepción
    llamaba por teléfono.
\end{itemize}

El punto de conflicto está en el dato \emph{cupos libres}. Ese valor solo
existe en Agenda, pero quien lo necesita es Reservas. Cualquier solución
tiene que resolver esa pregunta: ¿cómo se le pregunta a Agenda si hay
cupo, sin que los dos sistemas se mezclen?

\subsection{Qué se pidió}

El encargo fijaba un mínimo obligatorio:

\begin{enumerate}[leftmargin=1.7cm]
    \item Una \textbf{API REST pública} sobre dueños y reservas, que fuera
    el punto de entrada del sistema de Reservas hacia el exterior.
    \item Un \textbf{servicio interno en gRPC} (Agenda) que Reservas
    consumiera para verificar la disponibilidad de los veterinarios.
    \item La \textbf{integración} de ambos, de modo que una reserva solo se
    registrara si Agenda confirmaba que había cupo libre.
\end{enumerate}

Sobre Reservas se pedían las operaciones de administrar dueños (crear,
consultar y listar) y administrar reservas (reservar, revertir, consultar
y listar). Sobre Agenda se pedían tres operaciones: consultar un
veterinario y sus cupos, listar los bloques disponibles, y reservar o
liberar un cupo.

\subsection{Restricciones que hubo que respetar}

El equipo identificó cuatro restricciones que condicionaron todas las
decisiones posteriores:

\begin{enumerate}[leftmargin=1.7cm]
    \item \textbf{Cada sistema tiene un dueño distinto.} Agenda lo mantiene
    otro equipo, que sigue desarrollándolo de forma activa. Esto descarta
    cualquier opción en la que Reservas escriba en la base de Agenda, o en
    la que el equipo se apropie de ese código.
    \item \textbf{La consulta de disponibilidad es el camino más
    caliente.} Se pide "con muchísima frecuencia". El rendimiento importa
    más en ese camino que en el resto.
    \item \textbf{Todo debe funcionar con Docker.} No se acepta que
    funcione "en mi máquina". El sistema completo tiene que levantarse con
    un solo comando.
    \item \textbf{Quién nos consume no se actualiza a la vez.} Una
    colección de Postman, un script de la clínica o un portal futuro van a
    quedar congelados en la versión del día en que se escribieron. Por eso
    el contrato tiene que poder evolucionar sin romper a esos clientes.
\end{enumerate}

La cuarta restricción es la más fácil de pasar por alto y la que más
trabajo dio: no se trata solo de diseñar una API, sino de diseñar una
API que se pueda cambiar más adelante.

\subsection{Qué se decidió construir}

El equipo resolvió el problema de disponibilidad con tres decisiones
concretas, que se detallan en la sección \ref{sec:diseno} y se justifican
en la sección \ref{sec:adr}:

\begin{enumerate}[leftmargin=1.7cm]
    \item \textbf{No unir los sistemas en uno solo.} Se mantuvo Agenda y
    Reservas como servicios separados, cada uno con su propia base de
    datos. La comunicación entre ellos es un contrato gRPC, no un acceso
    directo a las tablas del otro.
    \item \textbf{Usar cada protocolo donde conviene.} La cara pública es
    REST, porque un navegador la puede consumir. La comunicación interna es
    gRPC, porque es tráfico entre dos servicios que se despliegan juntos y
    lo que importa es que sea barata y tipada.
    \item \textbf{Fallar de forma honesta cuando la dependencia no
    responde.} Si Agenda no contesta, la API responde 503 y no crea la
    reserva, en vez de inventar una disponibilidad.
\end{enumerate}

La tercera decisión es la que evita el problema original. Si la API
devolviera "hay cupo" usando datos de su propia caché cuando Agenda está
caído, volvería a existir la sobreventa. Preferimos fallar.

\subsection{Requisitos técnicos cubiertos}

La siguiente tabla muestra cómo se cubre cada requisito obligatorio del
encargo. Es la referencia que se usó durante el desarrollo para no dejar
ninguno atrás.

\begin{table}[H]
\centering
\caption{Requisitos técnicos obligatorios y su cobertura.}
\begin{tabular}{@{}p{1.1cm}p{4.2cm}p{8.2cm}@{}}
\toprule
\textbf{Req.} & \textbf{Descripción} & \textbf{Dónde se resuelve} \\
\midrule
T1 & Todo dockerizado &
Cada servicio es un contenedor. \texttt{docker-compose.yml} define los
tres (agenda, reservas, redis\_cache) y los levanta con un solo comando. \\
\addlinespace
T2 & API REST versionada &
Rutas bajo \texttt{/v1}, con verbos y códigos de estado correctos
(200, 201, 204, 400, 401, 404, 405, 409, 415, 503) y todos los errores en
JSON. \\
\addlinespace
T3 & Contrato explícito &
\texttt{contracts/openapi.yaml} para la API pública y
\texttt{contracts/agenda.proto} para el servicio interno. Ambos están en
el repositorio y versionados. \\
\addlinespace
T4 & Servicio interno en gRPC &
El servicio \texttt{AgendaService} del \texttt{.proto}, consumido por
Reservas mediante \texttt{agenda\_client.py}. Las tres operaciones son de
tipo Unary, y el ADR-002 explica por qué. \\
\addlinespace
T5 & Base de datos por servicio &
Cada servicio tiene su propio archivo SQLite, en su propio volumen de
Docker. Reservas nunca abre la base de Agenda. \\
\addlinespace
T6 & Autenticación en la API REST &
API Key en el header \texttt{X-API-Key}, validada por un decorador. El
ADR y la sección \ref{sec:contexto} justifican la elección. \\
\addlinespace
T7 & Manejo de fallas de la dependencia &
Timeout de 3 segundos en el cliente gRPC y respuesta 503 con JSON. Lo
detalla el ADR-004. \\
\bottomrule
\end{tabular}
\end{table}

% ==========================================
% 4. DISEÑO DE LA SOLUCION
% ==========================================

\newpage
\section{Diseño de la solución}\label{sec:diseno}

\subsection{Arquitectura general}

La figura \ref{fig:arquitectura} muestra el sistema completo. El
rectángulo punteado es la red de Docker. Lo único que sale de esa red
hacia el exterior es el puerto 5000 del servicio Reservas.

\begin{figure}[H]
\centering
\begin{tikzpicture}[
    font=\footnotesize,
    servicio/.style={draw, rounded corners=3pt, minimum width=3.4cm,
        minimum height=1.5cm, align=center, fill=blue!7, line width=0.6pt},
    almacen/.style={draw, dashed, rounded corners=3pt, minimum width=2.2cm,
        minimum height=1.1cm, align=center, fill=gray!10, line width=0.6pt},
    actor/.style={draw, rounded corners=3pt, minimum width=3.2cm,
        minimum height=0.9cm, align=center, fill=orange!12, line width=0.6pt},
    red/.style={draw, dashed, rounded corners=8pt, inner sep=10pt,
        fill=green!3, line width=0.7pt},
    flecha/.style={-{Stealth[length=2.2mm]}, thick},
    etq/.style={font=\scriptsize, text=black!65, midway}
]

% --- Actores externos ---
\node[actor] (navegador) at (0, 4.4) {Navegador\\ \texttt{localhost:5000}};
\node[actor] (postman) at (5.2, 4.4) {Postman / \texttt{curl}};

% --- Red de Docker ---
\node[red] (red) at (2.6, 1.1) {};

% --- Servicios ---
\node[servicio] (reservas) at (2.6, 1.9)
    {\textbf{Reservas}\\ API REST + interfaz web\\ \texttt{/v1} \quad puerto 5000};
\node[servicio] (agenda) at (-1.5, 0.2)
    {\textbf{Agenda}\\ servicio gRPC\\ puerto 50051};
\node[servicio] (redis) at (6.7, 0.2)
    {\textbf{Redis}\\ caché e idempotencia};

% --- Almacenes ---
\node[almacen] (bdres) at (2.6, -1.3) {Archivo SQLite\\\texttt{reservas.db}};
\node[almacen] (bdag) at (-1.5, -1.3) {Archivo SQLite\\\texttt{agenda.db}};

% --- Conexiones ---
\draw[flecha] (navegador) -- node[etq, above] {HTTP / REST} (reservas);
\draw[flecha] (postman) -- node[etq, above] {HTTP / REST} (reservas);
\draw[flecha] (reservas) -- node[etq, above] {\textbf{gRPC Unary}} (agenda);
\draw[flecha] (reservas) -- node[etq, above] {caché} (redis);
\draw[flecha] (agenda) -- (bdag);
\draw[flecha] (reservas) -- (bdres);

\end{tikzpicture}
\caption{Arquitectura del sistema VidaAnimal. Solo el servicio Reservas
publica un puerto hacia el host; Agenda y Redis son internos.}
\label{fig:arquitectura}
\end{figure}

\subsection{Descomposición en servicios}

El sistema quedó dividido en tres servicios, más un almacén de datos por
cada servicio de aplicación.

\begin{table}[H]
\centering
\caption{Los servicios del sistema y su responsabilidad.}
\begin{tabular}{@{}p{2.6cm}p{5.6cm}p{4.3cm}@{}}
\toprule
\textbf{Servicio} & \textbf{Responsabilidad} & \textbf{Tecnología} \\
\midrule
Agenda &
Administra los veterinarios y los cupos libres de cada bloque horario.
Es la fuente de verdad del cupo. &
Python, gRPC, SQLite \\
\addlinespace
Reservas &
Administra los dueños y las reservas. Es la puerta de entrada pública:
verifica el cupo contra Agenda antes de confirmar una cita. &
Python, Flask, SQLite \\
\addlinespace
Redis &
Guarda la lista de bloques durante 60 segundos y recuerda las respuestas
de las reservas idempotentes. &
Redis 7 \\
\bottomrule
\end{tabular}
\end{table}

\subsection{Justificación de la frontera}

La frontera entre Agenda y Reservas no es arbitraria, y esa es la decisión
más cara de todo el proyecto. El equipo la justificó con tres razones:

\begin{enumerate}[leftmargin=1.7cm]
    \item \textbf{Los dos sistemas responden preguntas distintas.} Agenda
    responde "¿hay cupo en este bloque?". Reservas responde "¿quién es
    este dueño y qué citas tiene?". Son dos contextos delimitados
    distintos, con datos que no se mezclan.
    \item \textbf{Cada uno tiene un equipo dueño distinto.} Como Agenda lo
    mantiene otro equipo, la frontera respetuosa es no tocar su base. Un
    servicio que escribe en la base de otro deja de ser autónomo: cada
    cambio de esquema rompe a un desconocido.
    \item \textbf{La frecuencia de uso es distinta.} La disponibilidad se
    consulta muchísimo más que lo que se reserva. Esa diferencia justifica
    una dependencia declarada y vigilada (con \emph{timeout} y caché), en
    vez de una llamada implícita a una tabla ajena.
\end{enumerate}

La consecuencia práctica es que la dependencia queda \emph{visible} en el
código: hay un cliente gRPC explícito (\texttt{agenda\_client.py}) que se
puede leer, en vez de un \texttt{SELECT} contra una base que no le
pertenece.

\subsection{El servicio interno Agenda (gRPC)}

El contrato está en \texttt{contracts/agenda.proto} y declara tres
operaciones, todas de tipo \textbf{Unary} (un mensaje de ida y uno de
vuelta). La tabla \ref{tab:grpc} las resume.

\begin{table}[H]
\centering
\caption{Operaciones del servicio gRPC Agenda.}
\label{tab:grpc}
\begin{tabular}{@{}p{3.1cm}p{3.6cm}p{6.0cm}@{}}
\toprule
\textbf{Operación} & \textbf{Tipo} & \textbf{Qué hace} \\
\midrule
\texttt{ListarBloques} & Unary &
Devuelve los bloques horarios disponibles, con filtros opcionales por
veterinario y por fecha. \\
\addlinespace
\texttt{ReservarCupo} & Unary &
Resta 1 al cupo libre del bloque. Si ya no quedaban cupos, responde con
un fallo y un mensaje. \\
\addlinespace
\texttt{LiberarCupo} & Unary &
Suma 1 al cupo libre del bloque. Se llama cuando se cancela una reserva. \\
\bottomrule
\end{tabular}
\end{table}

El descuento del cupo no es un \texttt{UPDATE} común: usa la condición
\texttt{WHERE cupos\_libres > 0}. Así, si dos peticiones llegan al mismo
tiempo, el segundo \texttt{UPDATE} no encuentra filas y la respuesta indica
que no había cupo. Es una forma barata de evitar la sobreventa dentro de
un solo proceso.

\subsection{La API REST de Reservas}

La API pública está bajo \texttt{/v1} y usa API Key. La tabla
\ref{tab:rest} lista las operaciones que quedaron implementadas.

\begin{table}[H]
\centering
\caption{Endpoints de la API REST de Reservas.}
\label{tab:rest}
\begin{tabular}{@{}p{1.2cm}p{4.4cm}p{1.6cm}p{5.5cm}@{}}
\toprule
\textbf{Verbo} & \textbf{Ruta} & \textbf{Código} & \textbf{Propósito} \\
\midrule
POST & \texttt{/v1/duenos} & 201 & Registra un dueño de mascota. \\
GET & \texttt{/v1/duenos} & 200 & Lista los dueños. Acepta filtro por nombre. \\
GET & \texttt{/v1/duenos/\{id\}} & 200 & Devuelve un dueño por su id. \\
\addlinespace
GET & \texttt{/v1/bloques} & 200 & Lista los bloques libres. Vienen de Agenda
por gRPC y se cachean 60 segundos. \\
\addlinespace
POST & \texttt{/v1/reservas} & 201 & Crea una reserva. Primero asegura el cupo
en Agenda. \\
GET & \texttt{/v1/reservas} & 200 & Lista las reservas. Acepta filtro por dueño
y por estado. \\
GET & \texttt{/v1/reservas/\{id\}} & 200 & Devuelve una reserva con enlaces
HATEOAS. \\
DELETE & \texttt{/v1/reservas/\{id\}} & 204 & Cancela la reserva y libera el cupo
en Agenda. \\
\addlinespace
GET & \texttt{/v1/health} & 200 & Indica que el servicio está arriba. No pide
API Key. \\
\bottomrule
\end{tabular}
\end{table}

Todos los errores usan la misma forma, con un código y un detalle:

\begin{lstlisting}[caption={Formato único de error.}]
{"error": "NO_ENCONTRADO", "detalle": "Recurso no encontrado"}
\end{lstlisting}

Cada código HTTP se escribe una sola vez, en su manejador global de
errores. Las rutas no arman el JSON: llaman a \texttt{abort(codigo)}, que
lanza el error y Flask encamina la respuesta por el manejador. Así un
dueño que no existe, una reserva que no existe y una URL mal escrita
devuelven exactamente el mismo cuerpo.

\subsection{Los contratos}

El sistema tiene dos contratos, y los dos están versionados en el
repositorio:

\begin{itemize}[leftmargin=1.4cm]
    \item \texttt{contracts/openapi.yaml}: describe la API REST pública.
    Se publica en \texttt{/v1/docs} con Swagger UI, que lee el mismo
    archivo. No hay una segunda copia de la especificación.
    \item \texttt{contracts/agenda.proto}: describe el servicio interno.
    El Dockerfile genera los stubs de Python al construir la imagen, así
    que si alguien agrega un campo al \texttt{.proto}, el compilador
    obliga a los dos lados a actualizarse.
\end{itemize}

\subsection{Almacenamiento de datos}

Cada servicio usa un archivo SQLite, guardado en su propio volumen de
Docker. SQLite no tiene servidor ni usuario ni contraseña: es un archivo.
La ventaja es que no hay que levantar una base de datos aparte, lo que
cumple con la exigencia de que todo funcione con un solo comando. La
desventaja es que no sirve para varias instancias del mismo servicio
escribiendo a la vez, algo que el equipo asume porque el alcance es de un
MVP.

% ==========================================
% 5. LOS ADR
% ==========================================

\newpage
\section{Decisiones de diseño (ADR)}\label{sec:adr}

Cada decisión importante se documenta como un ADR (*Architecture
Decision Record*). Un ADR registra el contexto, las alternativas que se
consideraron, la decisión, el costo que se aceptó y las consecuencias. El
equipo eligió documentar cuatro, que son las que pide el encargo.

\subsection{ADR-001 \texorpdfstring{---}{---} Dos servicios con base propia, en vez de un monolito}

\textbf{Estado:} aceptada \hfill \textbf{Responde a:} D1, T5

\textbf{Contexto.} VidaAnimal tiene dos sistemas que nacieron separados y
nunca se hablaron, y cada uno tiene un equipo dueño distinto. Además, la
disponibilidad se consulta con mucha frecuencia, así que el camino de esa
consulta es el más caliente del sistema.

\textbf{Alternativas.} Se consideraron tres: un monolito con una sola base
y las tablas de ambos dominios juntos; un servicio de Reservas que
consulta directamente la base de Agenda; y dos servicios con base propia
comunicados por gRPC.

\textbf{Decisión.} Agenda y Reservas son dos servicios independientes,
cada uno con su base de datos, y la única forma de comunicarse es el
contrato gRPC.

\textbf{Justificación.} Es la única alternativa que mantiene un almacén por
servicio sin atar a los dos equipos. La opción de consultar la base del
otro sería más rápida de implementar, pero un cambio en la tabla de cupos
rompería a Reservas y dejarían de existir dos almacenes separados. El
monolito resolvería el problema de las transacciones, pero haría que el
equipo de Agenda pierda su sistema propio.

\textbf{Costo aceptado.} La tabla \texttt{reservas} guarda una copia de la
fecha, la hora y el veterinario, porque así las reservas se pueden listar
sin llamar a Agenda en cada consulta. El costo es que esa copia puede
quedarse desactualizada si alguien cambia un bloque directamente en
Agenda. Solo el cupo es verdad; la cita es una foto del momento de la
reserva. También se acepta que no hay transacción distribuida: si el
cupo se asegura en Agenda y después falla el registro de la reserva,
queda un cupo descontado sin reserva, y hoy esa corrección es manual (una
transacción compensatoria o patrón Saga quedó fuera del alcance).

\textbf{Consecuencias.} Cualquier cambio en la disponibilidad obliga a
tocar el \texttt{.proto} y a regenerar los stubs. Si el volumen de
reservas creciera, la escritura pasaría a golpear Agenda y habría que
mirar ahí el cuello de botella.

\subsection{ADR-002 \texorpdfstring{---}{---} REST hacia afuera, gRPC hacia adentro}

\textbf{Estado:} aceptada \hfill \textbf{Responde a:} D2, T2, T4

\textbf{Contexto.} La API de Reservas la consume el personal de la clínica
desde un navegador. La consulta a Agenda es interna, de alto volumen y no
la ve el exterior. Ese reparto del tráfico es lo que decide el protocolo
de cada tramo.

\textbf{Alternativas.} gRPC en ambos lados; REST en ambos lados; y REST
hacia afuera con gRPC hacia adentro.

\textbf{Decisión.} La API pública es REST y se documenta en
\texttt{openapi.yaml}. La comunicación con Agenda es gRPC sobre el
\texttt{.proto}, y las tres operaciones son Unary.

\textbf{Justificación.} REST afuera porque es lo que un navegador puede
consumir sin ayuda: la interfaz web del proyecto habla REST y funciona. Es
fácil de revisar con \texttt{curl} o con las herramientas del navegador, y
sus códigos de estado son parte del contrato. gRPC adentro porque es
tráfico entre dos servicios que se despliegan juntos, donde importa el
costo por mensaje y que el contrato no se rompa en silencio: el
\texttt{.proto} genera el cliente y los stubs, así que agregar un campo
obliga a los dos lados. El modo Unary encaja porque las tres operaciones
son consultas o actualizaciones puntuales; el streaming solo se
justificaría si el servidor empujara datos sin que el cliente pregunte, y
aquí siempre es el cliente quien pregunta.

\textbf{Costo aceptado.} Dos contratos que se pueden desincronizar, porque
un cliente generado desde el \texttt{.proto} no valida el
\texttt{openapi.yaml} ni al revés. Además, los errores de gRPC no se
traducen solos: cada error hay que mapearlo a un código HTTP a mano.

\textbf{Consecuencias.} La interfaz web y la colección de Postman solo
necesitan hablar HTTP. Cada vez que se toca el \texttt{.proto} hay que
regenerar los stubs y volver a construir la imagen.

\subsection{ADR-003 \texorpdfstring{---}{---} Versionado en la ruta y cómo evoluciona el contrato}

\textbf{Estado:} aceptada \hfill \textbf{Responde a:} D3, T3

\textbf{Contexto.} Quien consume los contratos no se actualiza al mismo
tiempo que el equipo. La colección de Postman o un portal futuro quedan
congelados en la versión del día en que se escribieron. Por eso un cambio
de contrato no se puede hacer "en el sitio".

\textbf{Alternativas.} No versionar y cambiar el contrato directamente;
versionar por header (\texttt{Accept}); o versionar en la ruta.

\textbf{Decisión.} La API pública vive bajo \texttt{/v1} y la parte interna
se versiona en el \texttt{package} del \texttt{.proto}
(\texttt{agenda.v1}). Un cambio incompatible se agrega como versión nueva,
nunca se edita la vieja.

\textbf{Justificación.} La versión queda a la vista en la URL, en los logs
del servidor y en el \texttt{curl} de quien prueba la API, sin tener que
mirar headers escondidos. Permite además una transición real: publicar
\texttt{/v2} mientras \texttt{/v1} sigue sirviendo a los clientes que no
se han actualizado.

El criterio que sigue el equipo para clasificar un cambio es el que
resume la tabla \ref{tab:compat}.

\begin{table}[H]
\centering
\caption{Qué cambios del contrato son compatibles y cuáles no.}
\label{tab:compat}
\begin{tabular}{@{}p{8.0cm}p{2.2cm}p{2.9cm}@{}}
\toprule
\textbf{Cambio} & \textbf{Compatible} & \textbf{Por qué} \\
\midrule
Agregar un campo opcional a la respuesta & Sí &
Los clientes viejos ignoran lo que no conocen. \\
\addlinespace
Agregar un campo opcional al cuerpo & Sí, con cuidado &
El servidor debe seguir aceptando cuerpos que no lo manden. \\
\addlinespace
Agregar un campo obligatorio & \textbf{No} &
Los clientes antiguos no lo mandan y el servidor empezaría a rechazar
peticiones que antes funcionaban. Obliga a \texttt{/v2}. \\
\addlinespace
Cambiar el tipo de un campo & \textbf{No} &
Los clientes lo interpretan mal. Obliga a \texttt{/v2}. \\
\addlinespace
Renombrar o quitar un campo & \textbf{No} &
El cliente viejo sigue mandándolo o leyéndolo. \\
\addlinespace
Agregar un código de error nuevo & Sí &
Solo hay que documentarlo. \\
\bottomrule
\end{tabular}
\end{table}

En resumen: \emph{agregar es compatible, quitar o endurecer no}.

\textbf{Costo aceptado.} Mantener dos versiones en paralelo cuesta el doble
de pruebas y el doble de documentación, por eso solo se abre \texttt{/v2}
cuando el cambio es realmente incompatible.

\textbf{Consecuencias.} Antes de tocar un contrato hay que preguntarse si
el cambio es compatible; si no lo es, el costo sube a duplicar la ruta y
migrar. Los errores documentados son parte del contrato: un 404 significa
"el recurso no existe" para cualquier recurso, y ese texto general es lo
que permite reutilizar un único esquema \texttt{Error}.

\subsection{ADR-004 \texorpdfstring{---}{---} Qué responde la API cuando Agenda no está}

\textbf{Estado:} aceptada \hfill \textbf{Responde a:} D4, T7

\textbf{Contexto.} Ninguna reserva se registra sin que Agenda confirme que
hay cupo. Agenda se puede apagar, reiniciar o quedar lento. La API es la
puerta de entrada, así que es ella la que decide qué le contesta al
usuario cuando la dependencia no responde.

\textbf{Alternativas.} Se consideraron cuatro: esperar sin límite;
cortar con un timeout corto y responder 503; responder "hay cupo" desde la
caché y seguir como si nada; o cortar con timeout y seguir validando por
cuenta propia.

\textbf{Decisión.} Timeout de 3 segundos en cada llamada gRPC. Si se supera,
el cliente lanza un error propio y la API responde:

\begin{lstlisting}[caption={Respuesta cuando Agenda no responde.}]
HTTP 503 Service Unavailable
{"error": "AGENDA_NO_DISPONIBLE",
 "detalle": "El servicio de Agenda no está disponible.
            Intenta de nuevo en un momento."}
\end{lstlisting}

Además, la reserva \emph{no se inserta}: el cupo se asegura en Agenda antes
de escribir en la base de Reservas, así que un fallo de Agenda no deja
reservas fantasma.

\textbf{Justificación.} Un 503 dice "el problema no soy yo, reintenta", y
como la creación de reservas es idempotente, el reintento del cliente no
duplica nada. Un 500 diría "hay un error mío" y el usuario no debería
reintentar. La diferencia no es cosmética: es la que le permite a un
cliente decidir qué hacer. La alternativa de responder "hay cupo" desde la
caché se descartó porque es mentirle al usuario, y el sobrecupo es
justamente el problema que el sistema viene a resolver.

\textbf{Costo aceptado.} El usuario ve un error donde esperaba una reserva,
que es la consecuencia directa de no mentirle. Además, cada petición a
una dependencia caída espera hasta el timeout: no hay circuit breaker, así
que si llegan muchas peticiones juntas con Agenda caído, cada una paga su
espera. Un \textit{circuit breaker} o un reintento (\textit{Retry}) con 
espera progresiva (\textit{Exponential Backoff}) sería el siguiente paso, 
y no está implementado.

\textbf{Consecuencias.} La caché mitiga parte del problema: \texttt{GET
/v1/bloques} se cachea 60 segundos, así que mientras Agenda esté caído las
lecturas ya cacheadas se siguen respondiendo. El sistema no se cae entero,
solo deja de poder refrescar la información. La caché se borra al reservar
y al liberar, así que nunca muestra un cupo que ya se usó.

% ==========================================
% 6. EXPERIMENTO (4 PAGINAS RESERVADAS)
% ==========================================

\newpage
\section{Experimento: latencia y tamaño de los mensajes}\label{sec:experimento}

\begin{aviso}
Esta sección está reservada para la experimentación de la Competencia 6.
Su contenido lo completa el equipo: la hipótesis, el método, los datos
medidos, el análisis y las conclusiones. Las cuatro páginas siguientes
están dejando listas para que el equipo escriba los resultados sin tener
que modificar el formato. La idea es dejar el espacio listo para
que el equipo escriba los resultados.
\end{aviso}

\subsection{Hipótesis}

\subsection{Método y materiales}

\subsection{Datos}

\newpage
\subsection{Análisis}

\subsection{Conclusiones y límites}

% ==========================================
% 7. COMPETENCIA 2: FACTOR DE CONTEXTO
% ==========================================

\newpage
\section{Consideraciones de contexto}\label{sec:contexto}

\subsection{El factor elegido: privacidad de los datos de los dueños}

De los factores posibles (seguridad de la información, costo de
infraestructura, accesibilidad, impacto operacional), el equipo eligió
tratar la \textbf{privacidad de los datos de los dueños de mascotas}. Son
datos personales: nombre, teléfono y correo electrónico. Una lista de dueños 
de una clínica veterinaria es información sobre personas que asistieron
a una consulta en un lugar de salud animal, y en muchos países esa
información tiene protección legal.

\subsection{Cómo influyó en el diseño}

La elección del factor de privacidad no fue decorativa: cambió tres cosas
concretas del diseño.

\begin{enumerate}[leftmargin=1.7cm]
    \item \textbf{Autenticación con API Key en todas las rutas.} Se
    descartó dejar la API abierta. El endpoint \texttt{/v1/health} es la
    única excepción, y no devuelve datos personales, solo el estado del
    servicio. La clave se valida en un decorador antes de ejecutar la
    ruta, así que un cliente sin clave no alcanza ni a entrar al código de
    la consulta.
    \item \textbf{Sin datos de dueño en la ruta interna.} El tráfico entre
    Reservas y Agenda no incluye ningún dato personal: solo identificadores
    de bloque y conteos de cupos. La base de Agenda, que es la que no
    controla este equipo, nunca recibe nombres ni teléfonos.
    \item \textbf{Errores que no filtran información.} Los mensajes de
    error son generales a propósito. Un 404 dice "Recurso no encontrado" y
    no revela si el recurso no existe o si el cliente no tenía permiso
    para verlo. Un mensaje más específico ayudaría al usuario, pero
    también confirmaría la existencia de datos que quizás no le
    corresponden.
\end{enumerate}

\subsection{Otros factores considerados}

El equipo también evaluó otros dos factores y decidió no priorizarlos:

\begin{itemize}[leftmargin=1.4cm]
    \item \textbf{Costo de infraestructura.} Se optó por SQLite en
    archivos en lugar de un motor de bases de datos completo. Es la opción
    más barata en recursos y cumple con el alcance, pero no escala a
    varias instancias.
    \item \textbf{Accesibilidad.} La interfaz web usa HTML y Bootstrap sin
    dependencias de JavaScript complejas, y la API es REST puro. Es
    accesible desde cualquier navegador moderno sin instalar nada.
\end{itemize}

% ==========================================
% 8. REQUISITOS OPCIONALES
% ==========================================

\newpage
\section{Requisitos opcionales}

El encargo proponía cinco mejoras opcionales. El equipo implementó tres,
dejó una a medias y no hizo una. La tabla \ref{tab:opcionales} dice el
estado de cada una y dónde se puede ver en el código.

\begin{table}[H]
\centering
\caption{Estado de los requisitos opcionales.}
\label{tab:opcionales}
\begin{tabular}{@{}p{0.9cm}p{4.0cm}p{2.3cm}p{5.4cm}@{}}
\toprule
\textbf{\#} & \textbf{Mejora} & \textbf{Estado} & \textbf{Dónde está} \\
\midrule
O1 & Caché con Redis, con invalidación al reservar o liberar &
Implementada &
\texttt{reservas/app.py}, en las funciones \texttt{leer\_bloques\_cacheados},
\texttt{guardar\_bloques\_cacheados} y \texttt{olvidar\_bloques}. La caché
dura 60 segundos. \\
\addlinespace
O2 & Idempotencia con la cabecera \texttt{Idempotency-Key} &
Implementada &
\texttt{crear\_reserva()} lee la cabecera. Si la clave ya se usó, responde
200 con la reserva ya creada, en vez de crear una nueva. \\
\addlinespace
O3 & HATEOAS con enlaces según el estado &
Implementada &
\texttt{obtener\_reserva()} arma el campo \texttt{\_links} con
\texttt{self} y \texttt{dueno}, y agrega \texttt{cancelar} solo si la
reserva sigue activa. \\
\addlinespace
O4 & Pruebas de contrato que verifiquen el OpenAPI y el \texttt{.proto} &
Parcial &
La validez del OpenAPI se comprueba, y una prueba compara la
especificación servida con las rutas reales. No hay una suite que valide
el \texttt{.proto} ni que se ejecute en cada construcción. \\
\addlinespace
O5 & Segundo cliente gRPC en otro lenguaje &
No implementada &
Todo el consumo de Agenda es desde Python. Se descartó por alcance: el
\texttt{.proto} ya genera stubs para cualquier lenguaje. \\
\bottomrule
\end{tabular}
\end{table}

% ==========================================
% 9. INSTRUCCIONES DE EJECUCION
% ==========================================

\newpage
\section{Instrucciones de ejecución}

\subsection{Requisitos}

Solo hace falta Docker y Docker Compose. El proyecto no tiene modo de
ejecución local: no se instala nada en el computador y no se levanta una
base de datos aparte.

\subsection{Levantar el sistema completo}

\begin{lstlisting}
docker compose up --build
\end{lstlisting}

Ese es el único comando necesario. El sistema completo queda arriba:
el servicio de Agenda (gRPC), el de Reservas (API REST más interfaz web)
y la caché Redis.

\subsection{Accesos}

\begin{table}[H]
\centering
\caption{Qué se puede acceder desde el computador.}
\begin{tabular}{@{}p{4.8cm}p{7.7cm}@{}}
\toprule
\textbf{Qué} & \textbf{Dónde} \\
\midrule
Interfaz web y API REST & \texttt{http://localhost:5000} \\
Documentación (Swagger UI) & \texttt{http://localhost:5000/v1/docs} \\
Agenda (gRPC) & Solo dentro de la red de Docker \\
Redis (caché) & Solo dentro de la red de Docker \\
\bottomrule
\end{tabular}
\end{table}

Para probar los endpoints desde Swagger UI, hay que apretar el botón
\textbf{Authorize} y pegar la clave \texttt{clave-secreta-vet-2026}, que
se envía como el header \texttt{X-API-Key}.

\subsection{Probar la API con Postman}

El repositorio incluye una colección de Postman con 32 peticiones y 137
aserciones que cubren los diez estados HTTP de la API. Se importa con
\emph{Import \texorpdfstring{$\rightarrow$}{->} Link/Text/File}. Las
pruebas se pasan datos de una a otra, así que el orden importa, y hay que
correrlas en tres pasos:

\begin{enumerate}[leftmargin=1.7cm]
    \item Con Agenda arriba, correr las carpetas 0 a 5.
    \item Con \texttt{docker compose stop agenda}, poner la variable
    \texttt{agenda\_caida} en \texttt{si} y correr la carpeta 6, que
    comprueba los 503.
    \item Con \texttt{docker compose start agenda}, correr la carpeta 7,
    que cancela una reserva.
\end{enumerate}

\subsection{Comprobar el contrato OpenAPI}

El contrato se puede validar sin instalar nada en el computador, usando un
contenedor:

\begin{lstlisting}
docker run --rm -v "$PWD/contracts:/contrato:ro" python:3.12-slim \
  sh -c "pip install --quiet openapi-spec-validator && \
         openapi-spec-validator /contrato/openapi.yaml"
\end{lstlisting}

\subsection{Datos y reinicios}

Cada servicio crea su archivo SQLite al arrancar. Agenda carga datos de
ejemplo la primera vez, y solo si su base está vacía. Los archivos viven
en volúmenes de Docker, así que los datos sobreviven a \texttt{docker
compose down} y a \texttt{docker compose up --build}. Para volver a los
datos de ejemplo hay que usar \texttt{docker compose down -v}.

% ==========================================
% 10. REFLEXION FINAL
% ==========================================

\newpage
\section{Reflexión final}

\subsection{Qué hizo bien el equipo}

El equipo tomó buenas decisiones, y conviene ser honesto sobre cuáles
fueron y qué las motivó.

La primera fue partir el problema en dos servicios desde el principio,
en vez de empezar con un monolito y separar después. Resultó más fácil de
razonar: cada equipo podía trabajar sobre su propio sistema sin pisarse,
y la frontera quedó justificada por los datos y no por la estructura
del código.

La segunda fue escribir los contratos antes del código. Con el
\texttt{.proto} y el \texttt{openapi.yaml} listos, el equipo pudo
trabajar en paralelo y, sobre todo, pudo detectar temprano cuándo una
decisión de diseño no cerraba.

La tercera fue no inventar cosas. El equipo decidió no mentir cuando
Agenda no responde, y eso es más difícil que decir "hay cupo" desde la
caché. Es la decisión que más se valora en la evaluación del proyecto.

\subsection{Qué se haría distinto}

Con la experiencia de este trabajo, el equipo tomaría tres decisiones de
forma distinta:

\begin{enumerate}[leftmargin=1.7cm]
    \item \textbf{Medir antes de decidir.} Se eligieron la caché y el
    modo Unary de gRPC con buenos argumentos, pero sin cifras propias que
    los respaldaran. Con los números del experimento a la vista, la
    decisión habría sido más fácil de defender en la defensa.
    \item \textbf{Validar el \texttt{.proto} automáticamente.} El
    requisito opcional O4 quedó a medias. Con el \texttt{.proto}
    versionado, una prueba que verificara que el servicio cumple el
    contrato habría costado poco y habría detectado cambios accidentales.
    \item \textbf{Anticipar el \emph{circuit breaker} desde el diseño.} El
    ADR-004 lo reconoce como el siguiente paso natural, pero quedó como
    tarea futura. Con la estructura del proyecto ya montada, era el
    momento de implementarlo.
\end{enumerate}

\subsection{Qué quedó pendiente}

Quedaron pendientes tres cosas concretas: un \emph{circuit breaker} para
no acumular peticiones cuando Agenda está caído, una compensación
automática (patrón Saga) para el caso en que el cupo se asegura en Agenda pero
después falla el registro de la reserva, y un cliente gRPC en un segundo
lenguaje para demostrar la interoperabilidad del \texttt{.proto}.

\subsection{Conclusión}

El sistema resuelto cumple con los siete requisitos técnicos obligatorios
y con las tres mejoras opcionales de mayor valor. Más allá de eso, el
aprendizaje principal del trabajo fue que en un problema de integración
la dificultad casi nunca está en el código: está en decidir quién es
dueño de cada dato, qué pasa cuando una parte no está, y cómo se avisa a
quien nos consume cuando algo cambia. Esas tres decisiones son las que
quedaron documentadas en los ADR, y son las que el equipo defendería con
más confianza en la sustentación oral.

\end{document}
